\section{Related Work}

\subsection{Static Reflection Proposals and Tooling}
The evolution of compile-time metaprogramming in ISO C++ has progressed from template metaprogramming tricks and macro preprocessors to standardized \texttt{constexpr} facilities in C++14/17/20 \cite{p1240}. Proposals P2996 \cite{p2996} and P1967 \cite{p1967} represent the culmination of this effort, providing direct syntax for reflection and resource inclusion. CTRE \cite{ctre} established the precedent that an entire domain-specific grammar, in that case regular expressions, can be tokenized and compiled to typed C++ artifacts inside \texttt{constexpr} evaluation rather than through an external tool; \texttt{ctpb} applies the same principle to a substantially larger grammar, proto3, and adds type synthesis rather than code matching. Boost.PFR \cite{boost_pfr} predates P2996 and offers structural (member-order) reflection over existing aggregates, which is sufficient for serializing a struct that already exists but cannot synthesize a struct's members from an external schema the way \texttt{ctpb} does.

\subsection{Alternative Serialization Approaches}
\begin{itemize}
    \item \textbf{reflect-cpp \cite{reflect_cpp}:} Provides C++20 reflection-based serialization for JSON, YAML, and binary formats. However, \texttt{reflect-cpp} derives schemas from existing C++ structs rather than compiling external IDL schemas, lacking support for proto3 wire field tags, packed encoding, and reserved schema validation.
    \item \textbf{FlatBuffers \cite{flatbuffers} and Cap'n Proto \cite{capnproto}:} Focus on zero-copy in-memory layout but rely on standalone external compiler binaries similar to \texttt{protoc}.
    \item \textbf{libprotobuf-lite:} A stripped-down runtime library for Protocol Buffers that minimizes binary footprint by omitting runtime reflection descriptors. However, it still mandates build-time code generation via \texttt{protoc}.
    \item \textbf{nanopb \cite{nanopb} and protozero \cite{protozero}:} Lightweight C and C++ protobuf implementations aimed at embedded and streaming contexts respectively. Both still require a schema-compilation step (nanopb via its own \texttt{protoc} plugin, protozero by hand-written accessor code) rather than eliminating code generation outright.
    \item \textbf{upb \cite{upb}:} Google's own small-footprint C implementation of protobuf, used as the runtime core for several language bindings. Like \texttt{libprotobuf-lite}, it reduces runtime overhead but still depends on \texttt{protoc}-generated descriptors.
\end{itemize}

\texttt{ctpb} bridges this divide by delivering native IDL ingestion at compile time, uniting the schema authority of proto3 with zero-build-step header-only C++ integration.
